\section*{Chapter \ref{ch:b1:quantum-introduction} --- Introduction to Quantum Physics}

\begin{solution}{exo:b1:quantum-introduction:1}
$h\nu$: \qty{100}{MHz}: $6.6 \times 10^{-26}$~J $= \qty{4.1e-7}{eV}$; \qty{500}{nm}:
$1240/500 = \qty{2.5}{eV}$; \qty{0.10}{nm}: \qty{12.4}{keV}; \qty{1}{MeV}: $\lambda =
1240/10^6 = \qty{1.2e-3}{nm} = \qty{1.2}{pm}$. Laser: $1240/633 = \qty{1.96}{eV} =
\qty{3.1e-19}{J}$, so $10^{-3}/3.1 \times 10^{-19} = 3.2 \times 10^{15}$ photons per
second.
\end{solution}

\begin{solution}{exo:b1:quantum-introduction:2}
$\lambda_0 = 1240/2.3 = \qty{540}{nm}$. At \qty{400}{nm}, $h\nu = \qty{3.1}{eV}$:
$E_{k,\max} = \qty{0.8}{eV}$, $V_s = \qty{0.8}{V}$. Doubling the intensity
doubles the current and leaves $V_s$ unchanged. At \qty{600}{nm}
(\qty{2.07}{eV} $< W$): nothing, at any intensity.
\end{solution}

\begin{solution}{exo:b1:quantum-introduction:3}
Electron: $1.23/\sqrt{100} = \qty{0.12}{nm}$; at \qty{1}{keV}: \qty{0.039}{nm}.
Neutron: $E = \qty{4.0e-21}{J}$, $p = \sqrt{2m_nE} = \qty{3.7e-24}{kg.m/s}$,
$\lambda = h/p = \qty{0.18}{nm}$. Ball: $6.6 \times 10^{-34}/10^{-3} = \qty{7e-31}{m}$.
The first three have $\lambda$ of the order of the atomic spacing and are
diffracted (electron, neutron and X-ray crystallography); the ball never.
\end{solution}

\begin{solution}{exo:b1:quantum-introduction:4}
$E_1 = h^2/8m_eL^2 = 4.39 \times 10^{-67}/(8 \times 9.11 \times 10^{-31} \times 10^{-18}) =
\qty{6.0e-20}{J} = \qty{0.38}{eV}$, $E_2 = 4E_1 = \qty{1.5}{eV}$; photon
\qty{1.13}{eV}, $\lambda = 1240/1.13 = \qty{1.1}{\micro m}$ (near infrared). Proton
in \qty{5}{fm}: $4.39 \times 10^{-67}/(8 \times 1.67 \times 10^{-27} \times 25 \times 10^{-30}) =
\qty{1.3e-12}{J} = \qty{8}{MeV}$ --- the scale of nuclear energies, from
confinement alone.
\end{solution}

\begin{solution}{exo:b1:quantum-introduction:5}
Slope $(2.05 - 0.40)/(4.0 \times 10^{14}) = \qty{4.1e-15}{V.s}$, $h = e \times 4.1 \times
10^{-15} = \qty{6.6e-34}{J.s}$. $W = h\nu - eV_s = 2.48 - 0.40 = \qty{2.1}{eV}$,
threshold $\lambda_0 = 1240/2.1 = \qty{600}{nm}$: caesium.
\end{solution}

\begin{solution}{exo:b1:quantum-introduction:6}
$\Delta\lambda = \qty{2.43}{pm}$: $\lambda' = \qty{12.4}{pm}$. Before: $1240/0.010 =
\qty{124}{keV}$; after: $1240/0.0124 = \qty{100}{keV}$; the electron takes
\qty{24}{keV}. For light, $\Delta\lambda/\lambda = 2.4 \times 10^{-12}/5 \times 10^{-7} = 5 \times
10^{-6}$: unmeasurable in 1923 --- X-rays made the shift a $25\%$ effect.
\end{solution}

\begin{solution}{exo:b1:quantum-introduction:7}
$p = \sqrt{2 \times 9.11 \times 10^{-31} \times 5 \times 10^4 \times 1.6 \times 10^{-19}} =
\qty{1.2e-22}{kg.m/s}$, $\lambda = \qty{5.5}{pm}$ (relativity corrects it to
\qty{5.4}{pm}); \omterm{ex:b1:wave-propagation:young}{fringe spacing} $\lambda D/d = 5.5 \times 10^{-12} \times 0.35/2 \times 10^{-6}
= \qty{1}{\micro m}$: invisible to the \omterm{def:b1:optical-instruments:eye}{eye}, hence the magnifying \omterm{def:b1:mirrors-thin-lenses:lens}{lens}. Ten
electrons: ten dots, apparently at random; $10^5$: the fringes. Each
electron lands at one point; the \emph{probability} of landing follows
the two-slit wave --- a single electron's wave goes through both slits.
\end{solution}

\begin{solution}{exo:b1:quantum-introduction:8}
Atom: $\Delta p \geq \hbar/2\Delta x = \qty{5.3e-25}{kg.m/s}$, $E \sim \Delta p^2/2m_e =
\qty{1.5e-19}{J} \approx \qty{1}{eV}$: the right scale. Nucleus: $\Delta p \geq
\qty{1.1e-20}{kg.m/s}$; $\Delta p^2/2m_e$ would be \qty{400}{MeV}, far above
$m_ec^2$, so use $E \approx \Delta p\,c = \qty{3e-12}{J} = \qty{20}{MeV}$: no force
in the nucleus could hold an electron that energetic --- nuclei contain
no electrons (beta electrons are created when emitted). Grain: $\Delta v \geq
\hbar/2m\Delta x = 1.05 \times 10^{-34}/(2 \times 10^{-9} \times 10^{-6}) = \qty{5e-20}{m/s}$:
irrelevant.
\end{solution}

\begin{solution}{exo:b1:quantum-introduction:9}
(a) $m_ev^2/r = e^2/4\pi\varepsilon_0r^2$: $v = \sqrt{e^2/4\pi\varepsilon_0m_er}$, $E =
\tfrac12m_ev^2 - e^2/4\pi\varepsilon_0r = -e^2/8\pi\varepsilon_0r$. (b) $m_evr = n\hbar$ and
$m_ev^2r = e^2/4\pi\varepsilon_0$: divide the square of the first by the second:
$m_er = n^2\hbar^2 \cdot 4\pi\varepsilon_0/e^2$, $r_n = n^2a_0$, and $E_n = -e^2/8\pi\varepsilon_0n^2a_0
= E_1/n^2$ with $a_0 = \qty{52.9}{pm}$, $E_1 = \qty{-13.6}{eV}$. (c) $v_1 =
\hbar/m_ea_0 = \qty{2.19e6}{m/s}$, $v_1/c = 1/137$. (d) $2 \to 1$: $\tfrac34 \times
13.6 = \qty{10.2}{eV}$, \qty{122}{nm}; $3 \to 2$: $(\tfrac14 - \tfrac19) \times 13.6 =
\qty{1.89}{eV}$, \qty{656}{nm}. (e) $\lambda_n = h/m_ev_n = nh/m_ev_1 = 2\pi na_0$
(since $m_ev_1a_0 = \hbar$), and $2\pi r_n = 2\pi n^2a_0 = n\lambda_n$.
\end{solution}

\begin{solution}{exo:b1:quantum-introduction:10}
(a) Replace $e^2$ by $Ze^2$: $a_0 \to a_0/Z$, $E_1 \to Z^2E_1$. He$^+$:
\qty{54.4}{eV}. U$^{91+}$: $92^2 \times 13.6 = \qty{115}{keV}$ --- but $v_1 = Z\alpha c
= 0.67c$: the non-relativistic model is only indicative there. (b) $a_0/207
= \qty{256}{fm}$, $E_1 = 207 \times 13.6 = \qty{2.8}{keV}$; Lyman $\alpha$:
\qty{2.1}{keV}, $\lambda = \qty{0.59}{nm}$ (X-rays). The muon's orbit is only
$300$ proton radii: its levels shift measurably with the proton's size,
which is how the proton radius was re-measured (2010). (c) The two
particles orbit their common centre: the \omterm{prop:b1:systems-of-points:twobody}{reduced mass} $m_e/2$ halves
every energy: $E_1 = \qty{-6.8}{eV}$, Lyman $\alpha$ at \qty{5.1}{eV}, \qty{243}{nm}.
\end{solution}

\begin{solution}{exo:b1:quantum-introduction:11}
(a) $h^2/8m_e = \qty{6.02e-38}{J.m^2}$. $R = \qty{2}{nm}$: $E_e = 6.02 \times 10^{-38}/
(0.13 \times 4 \times 10^{-18}) = \qty{0.72}{eV}$, $E_h = \qty{0.21}{eV}$, total $1.74 +
0.93 = \qty{2.67}{eV}$, \qty{464}{nm}: blue. \qty{3}{nm}: confinement $\times4/9$,
\qty{2.15}{eV}, \qty{577}{nm}: yellow. \qty{4}{nm}: $\times1/4$, \qty{1.97}{eV},
\qty{629}{nm}: red. (b) \qty{520}{nm} is \qty{2.38}{eV}: confinement \qty{0.65}{eV}
$= 0.93 \times (2/R)^2$, $R = \qty{2.4}{nm}$. (c) No confinement: $E_g$ alone,
\qty{713}{nm}. (d) $\delta E = -2 \times 0.1 \times 0.41 = \qty{-0.08}{eV}$, $4\%$ of
\qty{2.15}{eV}: about \qty{20}{nm} of spread --- pure colours need
monodisperse batches.
\end{solution}

\begin{solution}{exo:b1:quantum-introduction:12}
(a) $E(\Delta x) = \hbar^2/8m\Delta x^2 + \tfrac12m\omega^2\Delta x^2$; zero derivative at
$\Delta x^2 = \hbar/2m\omega$, where both terms equal $\hbar\omega/4$: $E_{\min} =
\hbar\omega/2$. (b) $\tfrac12 \times 1.055 \times 10^{-34} \times 8.3 \times 10^{14} = \qty{4.4e-20}{J}
= \qty{0.27}{eV}$; $T \sim \hbar\omega/k_B = \qty{6300}{K}$ --- at room temperature
the vibration is frozen in its \omterm{prop:b1:quantum-introduction:box}{ground state}. (c) $\Delta p = \hbar/2\Delta x =
\qty{1.05e-24}{kg.m/s}$, $E = \Delta p^2/2m_{\mathrm{He}} = 1.1 \times 10^{-48}/1.33 \times
10^{-26} = \qty{8e-23}{J} = \qty{0.5}{meV}$: half the binding --- the atoms cannot
sit still enough to crystallize; helium stays liquid down to absolute
zero unless squeezed (\qty{25}{bar}). (d) $E_1 = \pi^2\hbar^2/2mL^2$ against
$\hbar^2/2mL^2$: ten times larger, as it must be --- the bound is a lower
limit, and the box state has $\Delta x \approx 0.18L$, not $L/2$.
\end{solution}

\begin{solution}{pb:b1:quantum-introduction:1}
\textbf{1.} Light frees electrons from the cathode with various
energies; held at $-V_s$, the anode repels them and only those with
$E_k > eV_s$ arrive. The current vanishes when $eV_s = E_{k,\max}$.

\textbf{2.} $eV_s = h\nu - W$: slope $h/e$, intercept $-W/e$; $V_s = 0$ at
the threshold $\nu_0 = W/h$.

\textbf{3.} Extreme points: $(2.36 - 0.30)/(5.0 \times 10^{14}) = \qty{4.12e-15}{V.s}$,
$h = 1.602 \times 10^{-19} \times 4.12 \times 10^{-15} = \qty{6.60e-34}{J.s}$ (a fit on
all five gives $6.62$): within $0.5\%$ of $6.626 \times 10^{-34}$.

\textbf{4.} $W = h\nu - eV_s = 2.47 - 0.30 = \qty{2.2}{eV}$; $\nu_0 = W/h =
\qty{5.3e14}{Hz}$, $\lambda_0 = 1240/2.2 = \qty{560}{nm}$. A red pointer
(\qty{650}{nm}, \qty{1.9}{eV}) extracts nothing.

\textbf{5.} $V_s$ unchanged, current doubled: twice the photons, each of
the same energy. A wave would give each electron more energy with more
intensity, and a threshold in intensity rather than in frequency.

\textbf{6.} $10^{-6} \times 10^{-20} = \qty{1e-26}{W}$ on the atom; \qty{2.2}{eV} $=
\qty{3.5e-19}{J}$ takes $3.5 \times 10^7$~s --- a year. Emission is observed
within nanoseconds: the energy arrives in whole quanta.

\textbf{7.} $m_ev^2/r = e^2/4\pi\varepsilon_0r^2$: $v = \sqrt{e^2/4\pi\varepsilon_0m_er}$; $E =
\tfrac12m_ev^2 - e^2/4\pi\varepsilon_0r = -e^2/8\pi\varepsilon_0r$.

\textbf{8.} $(m_evr)^2 = n^2\hbar^2$ and $m_ev^2r = e^2/4\pi\varepsilon_0$: $r_n = n^2 \cdot
4\pi\varepsilon_0\hbar^2/m_ee^2 = n^2a_0$, $a_0 = (1.055 \times 10^{-34})^2/(8.99 \times 10^9
\times 9.11 \times 10^{-31} \times 2.57 \times 10^{-38}) = \qty{52.9}{pm}$.

\textbf{9.} $E_n = -e^2/8\pi\varepsilon_0r_n = -m_ee^4/2(4\pi\varepsilon_0)^2\hbar^2n^2$; $E_1 =
-2.31 \times 10^{-28}/(2 \times 5.29 \times 10^{-11}) = \qty{-2.18e-18}{J} = \qty{-13.6}{eV}$.

\textbf{10.} $v_1 = \hbar/m_ea_0 = \qty{2.19e6}{m/s}$; $v_1/c = 7.3 \times 10^{-3} =
1/137$.

\textbf{11.} $2 \to 1$: \qty{10.2}{eV}, \qty{122}{nm} (ultraviolet); $3 \to 2$:
\qty{1.89}{eV}, \qty{656}{nm} (red); $4 \to 2$: \qty{2.55}{eV}, \qty{486}{nm}
(blue-green): the Balmer lines are the visible ones.

\textbf{12.} $hc/\lambda = |E_1|(1/n^2 - 1/p^2)$: $R_H = |E_1|/hc = 13.6/1240 =
\qty{1.097e-2}{nm^{-1}} = \qty{1.097e7}{m^{-1}}$. Balmer limit ($p \to \infty$):
$\lambda = 4/R_H = \qty{365}{nm}$.

\textbf{13.} \qty{13.6}{eV}; $\lambda = 1240/13.6 = \qty{91}{nm}$.

\textbf{14.} $v_n = v_1/n$, $\lambda_n = h/m_ev_n = nh/m_ev_1 = 2\pi na_0$ (as $m_ev_1a_0
= \hbar$); $2\pi r_n = 2\pi n^2a_0 = n\lambda_n$: the orbit carries $n$ whole
waves --- Bohr's rule is a standing-wave rule.

\textbf{15.} Wrong: the electron has no orbit (it would radiate; the
\omterm{prop:b1:quantum-introduction:box}{ground state} has zero \omterm{def:b1:angular-momentum:L}{angular momentum}) and the model fails for any atom
with two electrons. Right: the hydrogen levels, hence every line of its
spectrum. The orbit is replaced by a \omterm{def:b1:quantum-introduction:wavefunction}{wavefunction}, a probability cloud
of size $\sim n^2a_0$.

\textbf{16.} \omterm{thm:b1:wave-propagation:standing}{Standing waves}: $L = n\lambda/2$, $p_n = h/\lambda_n = nh/2L$, $E_n =
p_n^2/2m = n^2h^2/8mL^2$. The sphere formula is the same with $L \to R$:
the \omterm{prop:b1:quantum-introduction:box}{ground state} fits half a wavelength in the radius.

\textbf{17.} $E_e = 6.02 \times 10^{-38}/(0.13 \times 4 \times 10^{-18}) = \qty{0.72}{eV}$;
$E_h = 0.13/0.45 \times 0.72 = \qty{0.21}{eV}$.

\textbf{18.} $R = \qty{2}{nm}$: $1.74 + 0.93 = \qty{2.67}{eV}$, \qty{464}{nm}, blue;
\qty{3}{nm}: \qty{2.15}{eV}, \qty{577}{nm}, yellow; \qty{4}{nm}: \qty{1.97}{eV},
\qty{629}{nm}, red.

\textbf{19.} The confinement energy grows as $1/R^2$: a smaller box lifts
the levels, the photon is more energetic, the light bluer.

\textbf{20.} $\Delta p \sim \hbar/R = \qty{5.3e-26}{kg.m/s}$, $E \sim \Delta p^2/2m_e^\ast =
\qty{1.2e-20}{J} = \qty{0.07}{eV}$: the right order (the exact coefficient
$\pi^2/2$ and the cruder $\Delta p$ account for the factor ten).

\textbf{21.} $1240/577 = \qty{2.15}{eV} = \qty{3.4e-19}{J}$: $10^{-9}/3.4 \times 10^{-19}
= 2.9 \times 10^9$ photons per second.

\textbf{22.} $10^{-3}/(1.96 \times 1.6 \times 10^{-19}) = 3.2 \times 10^{15}$ per second.

\textbf{23.} $100 \times 2.48 \times 1.6 \times 10^{-19} = \qty{4e-17}{J}$ in \qty{0.1}{s}:
\qty{4e-16}{W}.

\textbf{24.} $h\nu = \qty{6.6e-26}{J}$: $1.5 \times 10^{13}$ photons per second ---
and each is $10^5$ times smaller than the thermal energy $k_BT$ of the
antenna: no hope, and no need; the wave picture is exact there.

\textbf{25.} The size of atoms, $a_0 = \qty{53}{pm}$; their energy scale,
\qty{13.6}{eV} (hence eV photons, chemistry, the visible spectrum); and
the speed of their electrons, $\alpha c \approx \qty{2e6}{m/s}$ --- all three
from $h$, $e$ and $m_e$ (and $\varepsilon_0$), with nothing adjusted.
\end{solution}
